\section{Conclusion}
\label{sec:conclusion}
In this paper, we revisit conventional single-modality conditioned 3D generation and highlight a clear limitation: image-conditioned models are sensitive to viewpoint informativeness and lack control over unobserved regions, while text-conditioned models capture global intent but lack concrete visual cues. Our diagnostic study empirically shows that these two signals are complementary, which motivates us to formalize the task of \textbf{Text--Image Conditioned 3D Generation}. We further introduce \textbf{TIGON}, a simple yet effective baseline for this task. Experiments demonstrate consistent gains over strong single-modality models, and show that combining modalities yields more robust and flexible 3D generation. We hope this work will help drive future research on controllable, high-quality 3D generation.



% In this paper, we revisit single-modality–conditioned 3D generation from the perspective of user control and find that it is hard to meet flexible specification needs: image-conditioned models are highly sensitive to viewpoint informativeness and struggle when the reference view underspecifies the underlying shape or appearance, while text-conditioned models capture global intent but lack concrete visual cues. Our diagnostic study empirically shows that these two signals are complementary, which motivates us to formalize the task of \textbf{Text--Image Conditioned 3D Generation}. We further introduce \textbf{TIGON}, a simple yet effective baseline for this task. Experiments demonstrate consistent gains over strong image-only and text-only models, and show that combining modalities yields more robust and flexible 3D generation. We hope this work will help drive future research on controllable, high-quality 3D generation.